\section{Methodology}

\subsection{Evaluation design}
The loop in \cref{fig:workflow} is replayed on historical data in a three-level, site-held-out design.
\begin{enumerate}
\item \emph{Source training.} In each of three folds the models are trained on source sites. Checkpoints
  and early stopping are chosen on that fold's validation sites.
\item \emph{Development.} Four replay sites select the primary model, the acquisition weights and the stop
  thresholds: Lower Hutt Town Hall, Tauranga Farmers, Christchurch Metro Sports and Marshland.
\item \emph{Evaluation.} Four replay sites test the frozen configuration: Wellington Airport, Christchurch
  Stadium, Hollywood and New Lock. Each is predicted with the fold weights that excluded it from both
  training and checkpoint selection.
\end{enumerate}

For every site and seed (17, 29, 43), whole groups are assigned to the audit set $A$ ($\approx$20\%), a
spatially dispersed $S_0$ ($\geq$5 positions) and the candidate pool $C$. Three policies share each split:
\begin{itemize}
\item \textbf{Full plan} (\Bfull) completes every position in $S_0\cup C$; this is the engineering reference.
\item \textbf{Sequential} (\Bseq) adds the two candidates that maximise the minimum distance to any
  revealed CPT. This \emph{maximin} rule is a space-filling baseline. It applies the stop rule.
\item \textbf{Adaptive} (\Padapt) adds the two candidates with the highest acquisition score and applies the
  same stop rule.
\end{itemize}
Full versus Adaptive measures the total benefit. Sequential versus Adaptive isolates what intelligent
location choice adds beyond simply stopping early. The evaluation therefore comprises 4~sites $\times$
3~seeds $\times$ 2~stopping policies $=24$ trajectories, plus 12 full-plan runs.

\textbf{Metrics.} All metrics are computed on $A$. (i)~\emph{log-MAE} is the mean $|\hat y-y|$ over audit cells; (ii)~\emph{thickness MAE} is the mean $|\hat H-H|$ over audit positions; and (iii)~the \emph{low-cell
metrics} are precision, recall, F1 and FPR.
Values are averaged within seed, then within site, and then equally across sites (the \emph{equal-site
mean}).

\subsection{Acquisition score}
For each candidate group the policy computes three signals: (i)~$U$, the mean GP predictive standard deviation of $\ln\qc$; (ii)~$D$, the share of cells where the primary model and XGBoost disagree on $\qc<1$\,MPa; (iii)~$d$, the distance to the nearest revealed CPT.
Each signal is converted to a 0--1 rank within the pool, and the score is
\begin{equation}
a(x)=0.5\,\mathrm{rank}(U)+0.3\,\mathrm{rank}(D)+0.2\,\mathrm{rank}(d).
\end{equation}
The weights are an engineering heuristic: uncertainty first, model disagreement second and spacing as a
tie-breaker. The only alternative tried omitted disagreement (0.7,~0,~0.3). Within a batch only $d$ is
updated after each pick, so the second choice is pushed away from the first without seeing its curve.

\subsection{Online stopping and independent audit}
Each round, the controller computes four signals on a fixed grid of all planned positions $\times$ 24
depths:
\begin{itemize}
\item the 90th percentile over grid cells of the round-to-round change in predicted $\ln\qc$ ($\Delta Q_{90}$);
\item the 90th percentile over positions of the change in predicted thickness ($\Delta H_{90}$);
\item a leave-group-out log-MAE on up to five initial groups;
\item the 90th percentile of the GP standard deviation.
\end{itemize}
The controller stops when all four pass for two consecutive rounds, provided at least nine CPTs are revealed
and every populated quadrant of the site is covered. Because batches add two CPTs, the earliest possible
stop is at 11~CPTs. Thresholds were set as quantiles (0.50/0.75/0.90) of the signals at acceptable
development states. The 0.50 quantile was selected because it maximised verified saving minus false-stop
rate; the 0.75 and 0.90 quantiles gave development false-stop rates of 17\% and 58\%.

The audit has two parts: (i)~\emph{absolute targets} of log-MAE $\le0.626$ and thickness MAE $\le0.881$\,m,
the development full-plan 75th percentiles plus fixed margins; and (ii)~a \emph{relative tolerance} of no
more than 0.05 log-MAE or 0.25\,m worse than the full plan.
Low-cell recall $\ge0.5$ with FPR $\le0.1$ is audited separately as a diagnostic.

\subsection{Hyperparameters}
\Cref{tab:hyper} lists every setting. Model settings were fixed during source training on validation sites.
Decision settings were fixed on development sites. No setting was tuned on evaluation sites.

\begin{table}[!htb]
\centering\small
\caption{Hyperparameters and how they were chosen.}
\label{tab:hyper}
\begin{tabularx}{\linewidth}{@{}lX@{}}
\toprule
Component & Setting (selection)\\
\midrule
GP & Mat\'ern $\nu=3/2$, anisotropic length scales initialised at (20, 20, 1)\,m, white-noise term; refitted
  by marginal likelihood each round\\
XGBoost residual & 3-NN IDW base (power 2) + $w\times$trees; 24 features; depth 6, learning rate 0.3,
  early stopping (36--37 trees); $w\in\{0, 0.25, 0.5, 1\}$ chosen per fold (0.25, 0.25, 0.5)\\
CNP & 64 hidden units, mean + max pooling, L1 loss on $\ln\qc$; 26,817 parameters; 500--1,200 steps,
  checkpoint by validation log-MAE\\
TabPFN & v2 weights (package 9.2.0), 1 estimator; fine-tuned with CRPS (continuous ranked probability score)
  + MSE, 500--600 updates; joint blend 0.5/0.5\\
Graph branch & 4 nearest neighbours, vertical weight 0.15, pseudo-Huber $\delta=0.35$\\
Acquisition & weights (0.5, 0.3, 0.2); batch of 2 groups; alternative (0.7, 0, 0.3) rejected on development\\
Stop rule & $\geq$9 CPTs, 2 rounds; $\Delta Q_{90}\le0.082$, $\Delta H_{90}\le0$\,m, support log-MAE
  $\le0.584$, GP-std P90 $\le0.461$ (0.50 quantile of development signals)\\
Experiment & seeds 17, 29, 43; audit share $\approx$20\%; $|S_0|\ge5$\\
\bottomrule
\end{tabularx}
\end{table}

\subsection{Economic model}
The owner's net benefit of a policy relative to the full plan is
\begin{equation}
B_{\text{net}}=(C_{\text{full}}-C_{\text{policy}})+\alpha\,\Delta T_{\text{inv}}\,(c_{\text{day}}+b_{\text{day}})-C_{\text{AI}},
\end{equation}
where the terms are: (i)~$C$ is investigation cost (\eur{750} per CPT, all-in); (ii)~$\Delta T_{\text{inv}}$ is the calendar days saved: 5 CPTs per rig-day with one rig, converted at 7/5, plus a 0.25-day decision delay at every batch transition; (iii)~$c_{\text{day}}=\eur{5,000}$ is the avoidable project cost per day; (iv)~$b_{\text{day}}=\eur{3,000}$ is the owner's benefit per day of earlier opening; (v)~$C_{\text{AI}}=\eur{600}$ is deployment and review per run.
In the main scenario $\alpha=1$, meaning every investigation day saved shortens the project. $\alpha=0.5$ and
0.75 are sensitivity cases, alongside one-at-a-time parameter ranges. A failed early stop would be repaired
to the full plan, with remobilisation (\eur{560}) and two days' delay.

Only the remobilisation cost and the consultant rate behind $C_{\text{AI}}$ come from a 2025 Dutch quotation
\parencite{mos2025quote}. All other values are explicit assumptions, following cost- and
schedule-estimating guidance \parencite{gao2015schedule,gao2020cost}. The same Dutch-based values are
applied to every site; for non-Dutch sites this is a standardised transfer, not local evidence.

\subsection{Tools and libraries}
The implementation uses Python 3.11, NumPy 2.4, pandas 3.0 and SciPy 1.17, scikit-learn 1.9 for the GP \parencite{pedregosa2011sklearn}, XGBoost 3.2, PyTorch 2.11 for the CNP, TabPFN 9.2.0 \parencite{priorlabs_tabpfn} and matplotlib.
Models ran on an NVIDIA RTX 4080 (16\,GB).
